\section{A Bespoke Special-Dividend Override}\label{sec:tour-special-div}

Real term sheets deviate from the standard treatment, and a system that hard-codes the
standard treatment has to be reprogrammed each time one does. This stop shows a
deviation absorbed as data. In September, ACME pays a special dividend, and the book
holds an option written over ACME whose confirmation adjusts the strike in a way the
class default would get wrong.

\subsection{Where the default and the term sheet disagree}

The option's strike is \EUR{100}. ACME declares a special dividend of \EUR{2} per
share, the boundary $(f, c) = (1, 2)$. The class default for a strike --- a
price-in-basis terms field --- under a special dividend is the full adjustment
$p \mapsto (p - c)/f$, which here is $K - c = 100 - 2 = \EUR{98}$. This default is
unconditional: once the event is classified special it applies in full, with no
threshold and no proportioning. And the classification itself --- ordinary or special
--- is the confirming authority's determination, recorded in the notice and made
outside the ledger; there is no ledger-side test that inspects the dividend's size and
decides the class. The ledger reads the class off the confirmed notice and applies the
default that class carries.

It is worth seeing why that default is what it is, because the invariance arithmetic
alone does not settle it. The re-expression $p \mapsto (p - c)/f$ preserves value only
when the same transaction pays the holder of the shifted field the matching cash leg
$q \cdot c$. A market quote crossing the ingest seam has that pairing --- the
shareholder who holds the quoted share is paid the dividend --- but the holder of an
option over the share is paid nothing. So Theorem~\ref{thm:dim-invariance} by itself
licenses no strike shift in either dividend class: the cash leg that would balance a
shift is simply not present for a strike. What decides the default is what the strike
already contains. An ordinary, anticipated dividend already sits in the forward the
strike was struck against, so the strike stands: shifting it down would hand the option
holder per share exactly the cash the boundary has already paid the shareholder --- one
distribution counted twice in the system's value accounting. A special dividend sits in
no forward, so a standing strike would impose that same unpriced transfer in the
opposite direction, this time against the holder; the strike moves to $K - c$ to hold
the holder's contracted position invariant. One principle yields both cases: a
price-in-basis terms field shifts by precisely the part of the distribution not already
embedded in the reference it was struck against --- zero for an ordinary dividend, the
whole of $c$ for a special. On an ordinary \EUR{2} dividend the strike therefore stands
at \EUR{100}, even as the read seam carries a \EUR{102} market print across the same
ex-boundary to \EUR{100}: the quote moves because its holder was paid, the strike does
not because the option holder was not. On the special \EUR{2} dividend the same
principle moves the strike the full $c$, to \EUR{98}.

That \EUR{98} is the default. This option's confirmation, though, says something
narrower: the strike adjusts on a special dividend only to the extent the dividend
exceeds \EUR{0.50}, and only by the excess. That is not the class default and not a
market norm --- it is this term sheet's own \emph{declared override}, written as a rule
in the schedule's closed first-order language that reads the boundary's declared cash
$c$ and the field's value $K$ and evaluates to
\[
  K \;\longmapsto\;
  \begin{cases}
    K - (c - 0.50) & \text{if } c > 0.50,\\
    K & \text{otherwise.}
  \end{cases}
\]
On this \EUR{2} dividend it yields $100 - (2 - 0.50) = \EUR{98.50}$ --- the excess-only
figure the confirmation requires, and the number the contract must produce. The rule is
total, and its boundary behaviour is exact: at $c = \EUR{0.50}$ the condition
$c > 0.50$ is false, so the strike stays \EUR{100}; one cent above, at $c = \EUR{0.51}$,
it moves by a single cent to $100 - 0.01 = \EUR{99.99}$. There is no gap and no
ambiguity at the threshold, because the rule evaluates to exactly one value for every
$c$. The threshold and the excess-only treatment live entirely in this override; the
class default behind it remains the full \EUR{98}.

\subsection{Three layers, and a gate that will not improvise}

The override is one of three ways a schedule entry can be filled, and the three form a
strict precedence, most specific winning:

\begin{enumerate}[leftmargin=1.5em]
\item the \textbf{class default}, fixed once in the specification per corporate-action
class and dimension --- what a strike does absent any deviation;
\item a \textbf{declared override}, a total rule in the closed language, reading the
boundary's parameters and the field's value --- where this term sheet's \EUR{98.50}
lives;
\item \textbf{designated discretion}, which names who owes a determination and leaves
the field with no adjusted value until that party's attested resolved terms arrive
referencing the boundary; until then the boundary is pending for that field and any
attempt to consume it is quarantined.
\end{enumerate}

What makes this safe rather than merely tidy is a totality gate. A unit is registered,
and every later terms version admitted, only if the schedule yields exactly one
well-formed entry for every pair of corporate-action class and terms field. The set of
classes is closed and finite, so the check is finite, and it is strict: an ill-formed
override is refused at the gate --- \texttt{ScheduleIncomplete} --- never quietly
defaulted to the class rule. The consequence is strong. Past the moment the option is
registered, ``the framework does not know how to adjust this field for this event'' is
not a state the system can reach. Either every class-and-field pair has exactly one
answer, or the version was never admitted.

So the special dividend runs through the same corporate-action contract as an ordinary
one. Nothing branches on the fact that this unit is bespoke. The contract reads the
schedule, finds the override entry, and appends the \EUR{98.50} strike; recomputation
from the logged notice re-runs the rule on $(K = 100,\ c = 2,\ \text{threshold } 0.50)$
and confirms \EUR{98.50}, diff zero. \emph{Bespokeness is a higher-precedence layer of
declared data behind a totality gate that fails closed --- never a special case in
code; the same contract that adjusts a plain strike adjusts a bespoke one, because the
deviation was written as data, not as a branch.}
